\documentclass[11pt,a4paper]{article}

\usepackage[margin=1in]{geometry}
\usepackage{booktabs}
\usepackage{xcolor}
\usepackage{hyperref}
\usepackage{amsmath}
\usepackage{enumitem}
\usepackage{fancyhdr}
\usepackage{titling}

\hypersetup{colorlinks=true, linkcolor=blue, urlcolor=blue}

\pagestyle{fancy}
\fancyhf{}
\fancyhead[L]{Bundle Density: Limits and Remaining Options}
\fancyhead[R]{\thepage}
\renewcommand{\headrulewidth}{0.4pt}

\setlength{\droptitle}{-1.5cm}
\title{\textbf{VLIW Bundle Density: Where the Limit Is, \\ and What Optimizations Remain}\\[2mm]
\large APARA C Compiler --- Optimization Status and Analysis}
\author{APARA C Compiler Project}
\date{\today}

\begin{document}
\maketitle
\thispagestyle{fancy}
\vspace{-6mm}

% ============================================================================
\section{Summary}
% ============================================================================

After five optimization passes (register caching, storage-class alias analysis, common
subexpression elimination, instruction scheduling, and loop-invariant code motion), the
compiler reaches an average of about \textbf{2.1 real instructions per bundle}, with the
parallel regions (vector ALU and dot-product kernels) already filling the \textbf{full
8-wide bundle}. This document explains, for the examination, \emph{why} the average is what
it is, whether it can be pushed substantially higher, and which optimization directions
remain worthwhile. The short conclusion: \textbf{bundle density is now close to the
architectural ceiling, and the productive direction for further speed-up is reducing the
number of executed loads/stores, not packing the bundles tighter.}

% ============================================================================
\section{What Determines Instructions-per-Bundle}
% ============================================================================

The number of instructions packed into a VLIW bundle is governed by three things, only one
of which the compiler controls:

\begin{enumerate}[leftmargin=*]
  \item \textbf{Available instruction-level parallelism (ILP)} --- how many mutually
        independent instructions exist to pack. This is a property of the program's data
        dependencies, not of the bundler.
  \item \textbf{The bundler's ability to exploit that ILP} --- our bundler is already a
        maximal greedy packer \emph{plus} a within-basic-block list scheduler that reorders
        independent instructions to cluster them. This part is effectively done.
  \item \textbf{The register file size} --- to issue $K$ independent operations in one
        cycle, their operands and results must all live in distinct registers
        simultaneously.
\end{enumerate}

The remaining headroom is limited by items (1) and (3), which are fundamental, not by (2).

% ============================================================================
\section{The Two Hard Limits}
% ============================================================================

\subsection*{Limit 1 --- the 28-register hardware ceiling}

APARA exposes 32 physical registers, of which 28 are allocatable (the rest are the zero
register, frame pointer, stack pointer, and global base). To pack $K$ \emph{independent}
arithmetic operations into one bundle requires roughly
\[
  \underbrace{2K}_{\text{operands}} \;+\; \underbrace{K}_{\text{results}} \;=\; 3K
  \quad\text{registers live at once.}
\]
With 28 registers, this caps the realistically packable independent operations at about
\textbf{8--9} --- which is exactly why the dense regions of our generated code already reach
the 8-instruction bundle width. \emph{One cannot pack more independent instructions than
there are registers to hold their operands.} This is a property of the hardware and cannot
be removed by the compiler.

\subsection*{Limit 2 --- genuine serial data dependencies}

The matrix-multiply kernel is an intrinsically serial chain:
\[
  \text{address} \;\to\; \text{wide-load} \;\to\; \texttt{\$dot} \;\to\;
  \texttt{\$dot \$accumulate} \;\to\; \text{store},
\]
where each step needs the previous step's result. No reordering or packing can parallelize
a \emph{single} instance of such a chain; the only parallelism comes from \emph{different,
independent} instances (e.g.\ separate loop iterations), and exploiting those requires more
registers (Limit~1) plus register renaming.

\medskip
\noindent Consequently the \emph{average} density (about 2.1) is held down by these serial
chains and by inherently-serial setup and control-flow code --- \textbf{not} by the bundler,
which is already maximal. The bundler is performing close to optimally given the register
file and the dependency structure.

% ============================================================================
\section{Remaining Density Techniques --- and an Honest Verdict}
% ============================================================================

\begin{table}[h]
\centering
\begin{tabular}{@{}p{4.3cm}p{4.3cm}p{4.6cm}@{}}
\toprule
Technique & Can it raise density? & Verdict \\
\midrule
Software pipelining & Yes --- overlap iteration $i{+}1$'s loads with iteration $i$'s
  compute (how the ISA's unrolled-dot example reaches 8-wide) & The legitimate way, but needs
  register renaming and is bounded by the 28-register wall; large, high-risk \\
Register renaming / SSA-based allocation & Enables unrolling and pipelining without false
  dependencies & Major allocator rewrite; \emph{still} capped at 28 registers, so limited
  headroom \\
Wider loads (\texttt{\$u256}) & Fewer load \emph{instructions} (more bytes per load) &
  Concrete but data-layout-dependent; modest gain \\
Instruction fusion / better selection & Fewer instructions overall & Marginal \\
\bottomrule
\end{tabular}
\end{table}

Every one of these runs into the same 28-register wall. That is the ``dead end'' one senses
when pushing density --- and it is real, but it is the \emph{hardware's} limit, not the
compiler's. The achievable gains are small relative to the implementation effort and risk
(both loop unrolling and the aggressive form of LICM already hit this register-pressure
ceiling during development).

% ============================================================================
\section{The Important Reframe: Density vs.\ Executed Work}
% ============================================================================

``More instructions per bundle'' (density) and ``faster execution'' (fewer executed
instructions) are different goals, and they diverged at the LICM stage:

\begin{itemize}[leftmargin=*]
  \item \textbf{Density} is near its architectural ceiling (limited by 28 registers and
        serial dependencies).
  \item \textbf{Executed work} is \emph{not} near a ceiling. Loop-invariant code motion cut
        the matrix multiply's executed loads by \textbf{33.5\%} (6441 $\to$ 4281) and total
        executed instructions by \textbf{22.7\%} (55201 $\to$ 42665) \emph{without changing
        bundle density at all} --- it simply does less work.
\end{itemize}

The lesson: for further speed-up, the productive direction is \textbf{doing less work}, not
packing bundles tighter. Reducing executed loads/stores improves runtime even when bundle
density stays the same.

% ============================================================================
\section{Recommended Next Steps (in priority order)}
% ============================================================================

\begin{enumerate}[leftmargin=*]
  \item \textbf{Loop-carried register allocation for loop counters} --- keep induction
        variables ($i$, $j$, $k$) in registers across iterations instead of reloading them
        from the stack every iteration. Pure executed-load reduction; same direction as
        LICM; does not fight the register wall. \emph{Highest value, lowest risk.}
  \item \textbf{Strength reduction} --- replace multiply/divide/modulo by a power of two
        with shift / shift / mask. Fewer ALU instructions; safe where signedness permits.
  \item \textbf{Wider loads} (\texttt{\$u256}) where the data layout allows --- fewer load
        instructions for the same bytes.
  \item \textbf{(Deferred / high-effort)} Software pipelining with register renaming ---
        the only path to materially higher \emph{density}, but bounded by the 28-register
        limit and a substantial, higher-risk rewrite. Recommended only if hot-loop
        throughput is the explicit objective.
\end{enumerate}

% ============================================================================
\section{Conclusion}
% ============================================================================

Bundle density is \textbf{not} blocked by a weak bundler; it is bounded by the APARA
register file (28 usable registers) and by genuine serial data dependencies. The dense
regions already reach the full 8-wide bundle, and the whole-program average of $\sim$2.1 is
close to optimal under those constraints. Pushing density materially higher would require
software pipelining plus register renaming --- a large effort with limited headroom because
of the 28-register ceiling. The better-value direction, already demonstrated by LICM's
22--33\% reduction in executed work, is to \textbf{continue eliminating redundant loads,
stores, and recomputation}, which improves runtime independently of bundle density.

\end{document}
